% R-06 authored by /root; complete argument checked and written by
% /root/fields_literature. Dependencies: finding environment, \BR,
% R-01--R-05, P-01, P-03, P-06 and the integral convention.
\section{The same minima for larger matrix sizes}

\begin{finding}{R-06}{uniform residual minima as the matrix size grows}
For both polynomial families, every algebraically closed field,
and every $n\geq4$,
\[
 \mu_3(U_{n,2})=8,\qquad
 \mu_4(U_{n,2})=9,\qquad
 \mu_5(U_{n,2})=11.
\]
In particular,
\[
 \BR(\mathcal D_n),\ \BR(\mathcal P_n)
 \ \geq\ \binom n2^{\!2}+6\qquad(n\geq4).
\]
For the two row families, every $n\geq2$ and
$1\leq j\leq n-1$ satisfy
\[
 \mu_j(U_{n,2})=j+1.
\]
All statements use the integral coefficient-tensor convention.
\end{finding}
\begin{proof}
P-01 reduces the polynomial assertion to sets of $j$
coordinate $2\times2$ rectangles. Let $E$ be the union of
their matrix-entry supports, and write $m=|E|$. As in the
finite certificates, the last-factor flattening has rank
$m$. Moreover, every used row and column contains at least
two elements of $E$, since each belongs to a selected
rectangle.

For $j=3$, the cases $m\geq8$ already have the required
lower bound. If $m\leq7$, at most three rows and three
columns occur, so row and column relabeling embeds the
configuration in a $4\times4$ matrix. R-02 applies there.
For $j=4$, the same argument with $m\leq8$ gives at most
four rows and four columns, so R-03 applies. Passing
between a tensor and its larger coordinate ambient space
does not change border rank: inclusion and coordinate
projection give the two inequalities.

For $j=5$ the cases $m\geq11$ are settled. If $m\leq10$,
at most five rows and five columns occur. A configuration
using at most four of each is handled by R-05. Suppose
that five rows occur. Then $m=10$, and every row has
exactly two column neighbors. Partition the five rows
according to their unordered pair of column neighbors.
Every part has at least two rows: a used row belongs to
a selected rectangle, and its other row has the same
two neighbors. A selected rectangle lies entirely in
one part. If the partition is $3+2$, at most
$\binom32+\binom22=4$ distinct rectangles exist, which
cannot contain the selected five. The only remaining
partition is $5$. Consequently the support is a
$5\times2$ grid. The case of five columns follows by
transposition.

The five selected rectangles in a $5\times2$ grid are
five edges on its five row vertices, with no isolated
vertex. There are exactly $222$ such labeled graphs:
from the $\binom{10}{5}=252$ five-edge subsets remove
the $5\binom65=30$ subsets with an isolated vertex.
Two isolated vertices are impossible for five edges.
The five permutation orbits and sufficient Koszul
certificates are as follows. Vertices in this table
are numbered $1,\ldots,5$; an entry $ij$ denotes an edge.
\[
\begin{array}{c|r|rr}
 \text{edge set}&\text{orbit size}&\Delta_{\det}&\Delta_{\mathrm{per}}\\ \hline
 12,13,14,15,23&30&-4&1\\
 12,13,14,23,25&60&-4&1\\
 12,13,14,23,45&60&4&1\\
 12,13,14,25,35&60&-1&1\\
 12,13,24,35,45&12&1&1
\end{array}
\]
Each entry is the determinant of an order-$41$ minor of
the $p=1$ Koszul map for the five-dimensional first
factor and the ten-dimensional support space. Its
pure-tensor cost is $\binom41=4$, so every nonzero
minor proves border rank at least
$\lceil41/4\rceil=11$. The certificate file
\texttt{five\_by\_two\_extra.json} specifies all row and
column indices. The separate checker
\texttt{verify\_five\_by\_two.py} reconstructs the maps
directly from the five ordered quadrics, computes these
ten determinants by integer elimination, and checks
that the disjoint permutation orbits cover all $222$
graphs. The permanent minors are units, and the
determinant minors are nonzero away from characteristic
two. In characteristic two the two tensors coincide.
This proves the lower bound in every characteristic.

The anchored shapes with $(h,r,s)=(3,1,3),(4,2,2),
(5,2,3)$ give upper bounds $8,9,11$ by P-06.
Each embeds in a $4\times4$ matrix, so these upper
bounds hold for all $n\geq4$. R-01 with $j=5$ now
gives $\binom n2^{2}-5+11=\binom n2^{2}+6$.

For a row residual space and any $j$-plane $W$, its
first factor is concise of dimension $j$. The all-field
one-dimensional residual theorem FL-03 gives matrix
rank at least two for every nonzero element of $W$.
Applying R-01 to $P_W$ while retaining one dimension
therefore gives $\BR(P_W)\geq j-1+2=j+1$.
For $j\leq n-1$, the coordinate star on the edges
$\{0,i\}$, $1\leq i\leq j$, is available. P-03 gives
its border rank $j+1$, proving equality.
\end{proof}

\begin{finding}{R-07}{monotonicity of the residual gain and the sixth minimum}
For any fixed residual space $U$ and
$1\leq i\leq j\leq\dim U$,
\[
 \mu_j(U)-j\ \geq\ \mu_i(U)-i.
\]
For both polynomial families in the integral convention,
over every algebraically closed field and every $n\geq4$,
\[
 \mu_6(U_{n,2})=12.
\]
\end{finding}
\begin{proof}
Apply R-01 to the concise first factor of $P_W$, where
$W\subset U$ has dimension $j$, retaining dimension $i$.
The surviving tensors are precisely $P_V$ for
$V\in\operatorname{Gr}(i,W)$. Hence
\[
 \BR(P_W)\geq j-i+
 \min_{V\in\operatorname{Gr}(i,W)}\BR(P_V)
 \geq j-i+\mu_i(U).
\]
Taking the minimum over $W$ proves the inequality.
R-06 then gives $\mu_6\geq6-5+11=12$.
The complete anchored $2\times3$ shape has six edges
and fits in a $3\times4$ matrix. P-06 gives its exact
border rank $6+2+3+1=12$, proving the matching upper
bound. This sixth minimum yields the same substitution
bound as the fifth minimum.
\end{proof}
